% Chapter 11: LLMs and Foundation Models for Time Series
% Harvard-quality academic presentation
% Bachelor program, Bucharest University of Economic Studies

\documentclass[9pt, aspectratio=169, t]{beamer}
\input{preamble}
\subtitle{Chapter 11: LLMs and Foundation Models for Time Series}

\begin{document}

% Title page (no header/footer)
{
\setbeamertemplate{headline}{}
\setbeamertemplate{footline}{}
\begin{frame}
    \titlepage
\end{frame}
}

%=============================================================================
% LEARNING OBJECTIVES
%=============================================================================
\begin{frame}{Learning Objectives}
    \begin{cminipage}{0.95\textwidth}
    \begin{block}{By the end of this chapter, you will be able to:}
        \begin{itemize}\setlength{\itemsep}{2pt}
            \item[\textcolor{MainBlue}{\textbf{1.}}] \textbf{Explain} the Transformer architecture and self-attention mechanism
            \item[\textcolor{MainBlue}{\textbf{2.}}] \textbf{Distinguish} encoder-only, decoder-only, and encoder-decoder architectures
            \item[\textcolor{MainBlue}{\textbf{3.}}] \textbf{Describe} the foundation model paradigm: pre-train, fine-tune, deploy
            \item[\textcolor{MainBlue}{\textbf{4.}}] \textbf{Compare} tokenization strategies: patching, quantization, and embeddings
            \item[\textcolor{MainBlue}{\textbf{5.}}] \textbf{Evaluate} major foundation models (Chronos, TimesFM, TimeGPT, Lag-Llama, Moirai)
            \item[\textcolor{MainBlue}{\textbf{6.}}] \textbf{Decide} when to use zero-shot inference vs.\ fine-tuning
            \item[\textcolor{MainBlue}{\textbf{7.}}] \textbf{Apply} foundation models to real-world forecasting tasks in Python
            \item[\textcolor{MainBlue}{\textbf{8.}}] \textbf{Assess} limitations, ethical concerns, and future directions
        \end{itemize}
    \end{block}
    \end{cminipage}
\end{frame}

%=============================================================================
% TABLE OF CONTENTS
%=============================================================================
\begin{frame}{Outline}
    \vspace{-0.3cm}
    {\small
    \setbeamertemplate{section in toc}{\color{MainBlue}$\boxdot$~\inserttocsection}
    \tableofcontents
    }
\end{frame}

%=============================================================================
\section{Motivation}
%=============================================================================

\begin{frame}{From Task-Specific to Foundation Models}
    \begin{cminipage}{0.95\textwidth}
    \vspace{-2mm}
    \begin{block}{Era 1: Classical Statistical Models (1970s--2010s)}
        \begin{itemize}\setlength{\itemsep}{0pt}
            \item ARIMA, ETS, GARCH --- hand-crafted per series, strong assumptions
            \item Strengths: interpretable, well-understood theory, works with small data
        \end{itemize}
    \end{block}
    \vspace{-2mm}
    \begin{alertblock}{Era 2: Deep Learning (2015--2022)}
        \begin{itemize}\setlength{\itemsep}{0pt}
            \item LSTM, TCN, DeepAR, N-BEATS --- trained on individual datasets
            \item Strengths: capture nonlinear dynamics; Weakness: require large labeled data per task
        \end{itemize}
    \end{alertblock}
    \vspace{-2mm}
    \begin{exampleblock}{Era 3: Foundation Models (2023--present)}
        \begin{itemize}\setlength{\itemsep}{0pt}
            \item TimeGPT, Chronos, TimesFM, Lag-Llama, Moirai --- pre-trained on billions of time points
            \item \textbf{Key insight}: train \textit{once} on massive diverse data, deploy \textit{everywhere} (zero-shot)
        \end{itemize}
    \end{exampleblock}
    \end{cminipage}
\end{frame}

\begin{frame}{The Scale Advantage}
    \begin{cminipage}{0.95\textwidth}
    \vspace{-2mm}
    {\footnotesize
    \begin{block}{Foundation Models: Scale Comparison}
    \vspace{-2mm}
    \begin{center}
    \begin{tabular}{l|r|l|l|c}
        \toprule
        \textbf{Model} & \textbf{Params} & \textbf{Pre-training Data} & \textbf{Architecture} & \textbf{Zero-Shot?} \\
        \midrule
        TimeGPT & N/A (proprietary) & 100B+ points & Transformer enc-dec & Yes (API) \\
        Chronos & 20M--710M & 84B points & T5 enc-dec & Yes \\
        TimesFM & 200M & 100B points & Decoder-only & Yes \\
        Lag-Llama & 352M & Multiple datasets & LLaMA decoder & Yes \\
        Moirai & 14M--311M & 27B points & BERT encoder & Yes \\
        PatchTST & 1M--10M & Task-specific & Transformer enc & No (supervised) \\
        \bottomrule
    \end{tabular}
    \end{center}
    \vspace{-2mm}
    \end{block}
    }
    \vspace{-2mm}
    \begin{exampleblock}{Key Observation}
        {\small Foundation models leverage \textbf{cross-domain knowledge transfer}: patterns learned from energy data help forecast retail sales, and vice versa. This is analogous to how GPT-3 learns grammar from diverse text corpora.}
    \end{exampleblock}
    \end{cminipage}
\end{frame}

\begin{frame}{Chapter Roadmap}
    \begin{cminipage}{0.95\textwidth}
    \vspace{-2mm}
    \begin{columns}[T]
        \column{0.55\textwidth}
        \begin{block}{What We Will Cover}
            {\small
            \begin{enumerate}\setlength{\itemsep}{1pt}
                \item \textbf{Motivation} --- the three eras of forecasting
                \item \textbf{Transformer Architecture} --- attention, positional encoding
                \item \textbf{From Transformers to Foundation Models} --- tokenization, patching
                \item \textbf{Deep Dive into Models} --- Chronos, TimesFM, TimeGPT, Lag-Llama, Moirai
                \item \textbf{Zero-Shot vs Fine-Tuning} --- when and how
                \item \textbf{Practical Use Cases} --- EUR/RON, energy, volatility
                \item \textbf{Benchmarks and Model Selection} --- Monash, GiftEval
                \item \textbf{Limitations, Ethics, and Future} --- what lies ahead
            \end{enumerate}
            }
        \end{block}

        \column{0.43\textwidth}
        \begin{center}
            \includegraphics[width=0.95\textwidth, height=0.50\textheight, keepaspectratio]{ch11_pretraining_scale.pdf}
        \end{center}
        {\scriptsize Pre-training data scale across foundation models (log scale).}
    \end{columns}
    \end{cminipage}
\end{frame}

%=============================================================================
\section{Transformer Architecture}
%=============================================================================

\begin{frame}{Researcher Spotlight: Vaswani et al.\ (2017)}
    \begin{cminipage}{0.95\textwidth}
    \vspace{-2mm}
    \begin{block}{``Attention Is All You Need''}
        {\small
        \begin{itemize}\setlength{\itemsep}{0pt}
            \item Published at NeurIPS 2017 by researchers at Google Brain
            \item Over \textbf{130{,}000 citations} --- one of the most influential ML papers ever
            \item Originally designed for \textbf{machine translation} (NLP)
            \item Key innovation: replaced recurrence entirely with \textbf{self-attention}
        \end{itemize}
        }
    \end{block}
    \vspace{-2mm}
    \begin{alertblock}{Why It Matters for Time Series}
        {\small
        \begin{itemize}\setlength{\itemsep}{0pt}
            \item Time series are \textbf{sequential} data --- same fundamental challenge as language
            \item Transformers handle \textbf{variable-length} sequences naturally
            \item \textbf{Parallel computation}: unlike RNNs, all positions processed simultaneously
            \item Foundation models (GPT, BERT, T5) all build on this architecture
        \end{itemize}
        }
    \end{alertblock}
    \vspace{-2mm}
    \begin{exampleblock}{Impact Timeline}
        {\small 2017: Attention Is All You Need $\to$ 2018: BERT, GPT $\to$ 2020: GPT-3 $\to$ 2023: TimeGPT, Chronos $\to$ 2024: TimesFM, Moirai}
    \end{exampleblock}
    \end{cminipage}
\end{frame}

\begin{frame}{Self-Attention: Intuition}
    \begin{cminipage}{0.95\textwidth}
    \vspace{-3mm}
    \begin{block}{Why Attention Replaces Recurrence}
        {\scriptsize
        \begin{itemize}\setlength{\itemsep}{0pt}
            \item \textbf{RNNs}: process tokens sequentially $\Rightarrow$ $O(T)$ sequential steps, information bottleneck
            \item \textbf{Self-Attention}: every position attends to every other position \textbf{in parallel}
            \item Direct connections between distant positions: no vanishing gradients
        \end{itemize}
        }
    \end{block}
    \vspace{-3mm}
    \begin{columns}[T]
        \column{0.48\textwidth}
        \begin{alertblock}{RNN: Sequential}
            {\scriptsize
            \begin{itemize}\setlength{\itemsep}{0pt}
                \item Path length between $x_1$ and $x_T$: $O(T)$ steps
                \item Cannot parallelize across time steps
                \item Gradient must flow through $T$ multiplications
                \item Information ``compresses'' through hidden state
            \end{itemize}
            }
        \end{alertblock}

        \column{0.48\textwidth}
        \begin{exampleblock}{Attention: Parallel}
            {\scriptsize
            \begin{itemize}\setlength{\itemsep}{0pt}
                \item Path length between any two positions: $O(1)$
                \item Fully parallelizable $\Rightarrow$ faster on GPUs
                \item Each position directly ``sees'' all others
                \item Learns which positions to attend to (data-driven)
            \end{itemize}
            }
        \end{exampleblock}
    \end{columns}
    \vspace{-2mm}
    \begin{center}
        \includegraphics[width=0.90\textwidth, height=0.22\textheight, keepaspectratio]{ch11_quiz1_attention_vs_rnn.pdf}
    \end{center}
    \end{cminipage}
\end{frame}

\begin{frame}{Self-Attention: Mathematics}
    \begin{cminipage}{0.95\textwidth}
    \vspace{-2mm}
    \begin{block}{Scaled Dot-Product Attention}
        Given input embeddings $X \in \mathbb{R}^{T \times d}$, compute:
        \vspace{-2mm}
        $$Q = XW^Q, \quad K = XW^K, \quad V = XW^V$$
        \vspace{-4mm}
        $$\text{Attention}(Q,K,V) = \text{softmax}\!\left(\frac{QK^\top}{\sqrt{d_k}}\right)V$$
    \end{block}
    \vspace{-2mm}
    \begin{columns}[T]
        \column{0.50\textwidth}
        {\scriptsize
        \begin{exampleblock}{Interpretation}
            \begin{itemize}\setlength{\itemsep}{0pt}
                \item $Q$ (Queries): ``What am I looking for?''
                \item $K$ (Keys): ``What do I contain?''
                \item $V$ (Values): ``What information do I provide?''
                \item $\sqrt{d_k}$: scaling prevents softmax saturation
                \item Output: weighted average of values, weights = compatibility scores
            \end{itemize}
        \end{exampleblock}
        }

        \column{0.48\textwidth}
        \begin{center}
            \includegraphics[width=0.95\textwidth, height=0.38\textheight, keepaspectratio]{ch11_attention_heatmap.pdf}
        \end{center}
        {\scriptsize Attention weights heatmap: rows attend to columns.}
    \end{columns}
    \end{cminipage}
\end{frame}

\begin{frame}{Multi-Head Attention}
    \begin{cminipage}{0.95\textwidth}
    \vspace{-2mm}
    \begin{block}{Multiple Attention Heads Capture Diverse Patterns}
        \vspace{-2mm}
        \begin{align*}
            \text{MultiHead}(Q,K,V) &= \text{Concat}(\text{head}_1, \ldots, \text{head}_h)\, W^O \\
            \text{where head}_i &= \text{Attention}(QW_i^Q,\; KW_i^K,\; VW_i^V)
        \end{align*}
        \vspace{-4mm}
        {\small Each head projects to a lower dimension: $W_i^Q, W_i^K \in \mathbb{R}^{d \times d_k}$, $W_i^V \in \mathbb{R}^{d \times d_v}$ with $d_k = d_v = d/h$.}
    \end{block}

    \vspace{-1mm}

    \begin{columns}[T]
        \column{0.48\textwidth}
        \begin{exampleblock}{What Different Heads Learn}
            {\scriptsize
            \begin{itemize}\setlength{\itemsep}{0pt}
                \item Head 1: \textbf{Trend} --- attend to recent values
                \item Head 2: \textbf{Seasonality} --- attend to same period last cycle
                \item Head 3: \textbf{Outliers} --- attend to extreme values
                \item Head 4: \textbf{Long-range} --- attend to distant correlations
            \end{itemize}
            }
        \end{exampleblock}

        \column{0.48\textwidth}
        \begin{alertblock}{Computational Cost}
            {\scriptsize
            \begin{itemize}\setlength{\itemsep}{0pt}
                \item Self-attention complexity: $O(T^2 \cdot d)$
                \item Memory: $O(T^2)$ for the attention matrix
                \item For long sequences ($T > 4096$): quadratic cost becomes prohibitive
                \item Solutions: patching, sparse attention, linear attention
            \end{itemize}
            }
        \end{alertblock}
    \end{columns}
    \end{cminipage}
\end{frame}

\begin{frame}{Positional Encoding}
    \begin{cminipage}{0.95\textwidth}
    \vspace{-2mm}
    \begin{block}{Why Position Matters}
        {\small Self-attention is \textbf{permutation-equivariant}: without positional info, it treats $\{x_1, x_2, x_3\}$ and $\{x_3, x_1, x_2\}$ identically. Time series \textbf{order matters}, so we inject position.}
    \end{block}
    \vspace{-2mm}
    \begin{alertblock}{Sinusoidal Positional Encoding (Vaswani et al., 2017)}
        \vspace{-2mm}
        \begin{align*}
            PE_{(pos,\, 2i)} &= \sin\!\left(\frac{pos}{10000^{2i/d}}\right) \\
            PE_{(pos,\, 2i+1)} &= \cos\!\left(\frac{pos}{10000^{2i/d}}\right)
        \end{align*}
        \vspace{-4mm}
        {\small Each dimension uses a different frequency $\Rightarrow$ unique positional signature. Relative positions can be computed via dot products.}
    \end{alertblock}
    \vspace{-2mm}
    \begin{center}
        \includegraphics[width=0.95\textwidth, height=0.28\textheight, keepaspectratio]{ch11_positional_encoding.pdf}
    \end{center}
    \end{cminipage}
\end{frame}

\begin{frame}{The Full Transformer Block}
    \begin{cminipage}{0.95\textwidth}
    \vspace{-2mm}
    \begin{columns}[T]
        \column{0.50\textwidth}
        \begin{block}{Transformer Block Components}
            {\scriptsize
            \begin{enumerate}\setlength{\itemsep}{0pt}
                \item \textbf{Input embedding} + Positional encoding
                \item \textbf{Layer Normalization} (pre-norm variant)
                \item \textbf{Multi-Head Self-Attention}
                \item \textbf{Residual connection}: $x + \text{Attn}(x)$
                \item \textbf{Layer Normalization}
                \item \textbf{Feed-Forward Network} (FFN): two linear layers with GELU/ReLU
                $$\text{FFN}(x) = W_2 \cdot \text{GELU}(W_1 x + b_1) + b_2$$
                \item \textbf{Residual connection}: $x + \text{FFN}(x)$
            \end{enumerate}
            }
        \end{block}
        {\scriptsize Stacking $N$ blocks ($N=6$--$24$) gives depth. Residual connections enable gradient flow.}

        \column{0.48\textwidth}
        \begin{center}
            \includegraphics[width=0.95\textwidth, height=0.55\textheight, keepaspectratio]{ch11_transformer_block.pdf}
        \end{center}
        {\scriptsize Diagram: the standard Transformer block with pre-norm layout.}
    \end{columns}
    \end{cminipage}
\end{frame}

\begin{frame}{Encoder vs Decoder Architectures}
    \begin{cminipage}{0.95\textwidth}
    \vspace{-2mm}
    {\footnotesize
    \begin{block}{Three Architectural Paradigms}
    \vspace{-2mm}
    \begin{center}
    \begin{tabular}{l|p{3cm}|p{3cm}|p{3.5cm}}
        \toprule
        & \textbf{Encoder-only} & \textbf{Decoder-only} & \textbf{Encoder-Decoder} \\
        \midrule
        \textbf{NLP origin} & BERT & GPT & T5, BART \\
        \textbf{Attention} & Bidirectional & Causal (masked) & Cross-attention \\
        \textbf{TS models} & Moirai, PatchTST & TimesFM, Lag-Llama & Chronos, TimeGPT \\
        \textbf{Strengths} & Learns rich representations & Natural for autoregressive generation & Flexible input-output mapping \\
        \textbf{TS use} & Representation, classification & Forecasting, generation & Forecasting, seq2seq \\
        \bottomrule
    \end{tabular}
    \end{center}
    \vspace{-2mm}
    \end{block}
    }
    \vspace{-2mm}
    \begin{exampleblock}{Key Distinction}
        {\small \textbf{Causal masking} (decoder): position $t$ can only attend to $\{1, \ldots, t\}$. Prevents ``peeking'' into the future --- natural fit for time series forecasting. \textbf{Bidirectional} (encoder): every position sees all others --- useful for representation learning and classification.}
    \end{exampleblock}
    \end{cminipage}
\end{frame}

%=============================================================================
\section{From Transformers to Foundation Models}
%=============================================================================

\begin{frame}{The Foundation Model Paradigm}
    \begin{cminipage}{0.95\textwidth}
    \vspace{-2mm}
    \begin{block}{Three-Stage Pipeline}
        {\small
        \begin{enumerate}\setlength{\itemsep}{0pt}
            \item \textbf{Pre-training}: Train on massive, diverse corpora (billions of time points across domains)
            \item \textbf{Fine-tuning} (optional): Adapt to a specific domain or dataset with a small learning rate
            \item \textbf{Inference}: Deploy on new, unseen series --- zero-shot or few-shot
        \end{enumerate}
        }
    \end{block}
    \vspace{-2mm}
    \begin{columns}[T]
        \column{0.50\textwidth}
        \begin{center}
            \includegraphics[width=0.95\textwidth, height=0.40\textheight, keepaspectratio]{ch11_foundation_paradigm.pdf}
        \end{center}

        \column{0.48\textwidth}
        \begin{alertblock}{Key Insight (Bommasani et al., 2021)}
            {\scriptsize
            \begin{itemize}\setlength{\itemsep}{0pt}
                \item ``A foundation model is any model trained on broad data that can be adapted to a wide range of downstream tasks.''
                \item For TS: train on \textit{energy + finance + retail + health + weather} $\Rightarrow$ forecast \textit{anything}
                \item Critical assumption: \textbf{shared temporal patterns} exist across domains (trend, seasonality, noise)
            \end{itemize}
            }
        \end{alertblock}
    \end{columns}
    \end{cminipage}
\end{frame}

\begin{frame}{Tokenization Strategies for Time Series}
    \begin{cminipage}{0.95\textwidth}
    \vspace{-2mm}
    {\scriptsize
    \begin{block}{Challenge: Continuous Values $\to$ Discrete Tokens}
        NLP models process discrete tokens (words). Time series are \textbf{continuous} $\Rightarrow$ need a mapping strategy.
    \end{block}
    \vspace{-2mm}
    \begin{columns}[T]
        \column{0.48\textwidth}
        \begin{exampleblock}{Strategy 1: Raw Value Embeddings}
            \begin{itemize}\setlength{\itemsep}{0pt}
                \item Project each $x_t \in \mathbb{R}$ to $\mathbb{R}^d$ via linear layer
                \item Used by: vanilla Transformer TS models
                \item Simple but ignores local structure
            \end{itemize}
        \end{exampleblock}
        \vspace{-2mm}
        \begin{exampleblock}{Strategy 2: Scaling + Quantization}
            \begin{itemize}\setlength{\itemsep}{0pt}
                \item Normalize, then map to $B$ discrete bins
                \item Used by: \textbf{Chronos} (4096 bins)
                \item Leverages NLP machinery directly
            \end{itemize}
        \end{exampleblock}

        \column{0.48\textwidth}
        \begin{exampleblock}{Strategy 3: Patching}
            \begin{itemize}\setlength{\itemsep}{0pt}
                \item Group $P$ consecutive points into one token
                \item Used by: \textbf{PatchTST}, \textbf{TimesFM}
                \item Reduces sequence length by factor $P$
            \end{itemize}
        \end{exampleblock}
        \vspace{-2mm}
        \begin{exampleblock}{Strategy 4: Variate-Independent Channels}
            \begin{itemize}\setlength{\itemsep}{0pt}
                \item Each variable processed independently
                \item Used by: \textbf{PatchTST}, \textbf{Moirai}
                \item Scales to any number of variates
            \end{itemize}
        \end{exampleblock}
    \end{columns}
    }
    \vspace{-1mm}
    \begin{center}
        \includegraphics[width=0.95\textwidth, height=0.18\textheight, keepaspectratio]{ch11_tokenization_strategies.pdf}
    \end{center}
    \end{cminipage}
\end{frame}

\begin{frame}{Patching: Why It Works}
    \begin{cminipage}{0.95\textwidth}
    \vspace{-2mm}
    \begin{block}{PatchTST Insight (Nie et al., 2023)}
        {\small Divide a time series of length $L$ into $N = \lfloor L / P \rfloor$ non-overlapping patches of size $P$:
        $$\mathbf{p}_i = (x_{(i-1)P+1},\; x_{(i-1)P+2},\; \ldots,\; x_{iP}) \in \mathbb{R}^P, \quad i = 1, \ldots, N$$
        Each patch is linearly projected to dimension $d$ and fed as a single token to the Transformer.}
    \end{block}
    \vspace{-2mm}
    \begin{columns}[T]
        \column{0.48\textwidth}
        \begin{exampleblock}{Benefits of Patching}
            {\scriptsize
            \begin{itemize}\setlength{\itemsep}{0pt}
                \item \textbf{Reduces complexity}: attention $O(T^2) \to O((T/P)^2)$
                \item \textbf{Local semantics}: each patch captures a short sub-pattern (e.g., daily cycle)
                \item \textbf{Longer context}: with $P=32$, a 512-token window covers 16{,}384 time steps
                \item \textbf{Better representations}: patch-level tokens are more informative than point-level
            \end{itemize}
            }
        \end{exampleblock}

        \column{0.48\textwidth}
        \begin{center}
            \includegraphics[width=0.95\textwidth, height=0.45\textheight, keepaspectratio]{ch11_patching_illustration.pdf}
        \end{center}
    \end{columns}
    \end{cminipage}
\end{frame}

\begin{frame}{Quantization: The Chronos Approach}
    \begin{cminipage}{0.95\textwidth}
    \vspace{-2mm}
    \begin{block}{Scale-and-Quantize Tokenization (Ansari et al., 2024)}
        {\small
        \begin{enumerate}\setlength{\itemsep}{0pt}
            \item \textbf{Normalize}: $\tilde{x}_t = \frac{x_t - \bar{x}}{s}$ where $\bar{x}$ is mean and $s$ is the mean absolute value
            \item \textbf{Quantize}: map to $B$ bins via the CDF: $b_t = \left\lfloor (B-1) \cdot \Phi(\tilde{x}_t) \right\rfloor$ where $\Phi$ is the standard normal CDF
            \item \textbf{Embed}: look up bin $b_t$ in the T5 embedding table ($B = 4096$)
        \end{enumerate}
        }
    \end{block}
    \vspace{-2mm}
    \begin{columns}[T]
        \column{0.48\textwidth}
        \begin{alertblock}{Why Quantize?}
            {\scriptsize
            \begin{itemize}\setlength{\itemsep}{0pt}
                \item Reuses \textbf{NLP vocabulary machinery}: embeddings, softmax output layer
                \item Cross-entropy loss (categorical) is well-understood
                \item Probabilistic: output is a distribution over bins $\Rightarrow$ full predictive distribution
                \item Trade-off: resolution limited by number of bins
            \end{itemize}
            }
        \end{alertblock}

        \column{0.48\textwidth}
        \begin{center}
            \includegraphics[width=0.95\textwidth, height=0.40\textheight, keepaspectratio]{ch11_chronos_tokenization.pdf}
        \end{center}
    \end{columns}
    \end{cminipage}
\end{frame}

\begin{frame}{Pre-Training Objectives}
    \begin{cminipage}{0.95\textwidth}
    \vspace{-3mm}
    {\scriptsize
    \begin{block}{Paradigm 1: Next-Token Prediction (Autoregressive, GPT-style)}
        \vspace{-1mm}
        \begin{itemize}\setlength{\itemsep}{0pt}
            \item Predict $x_{t+1}$ given $x_1, \ldots, x_t$ via causal masking. Loss: $\mathcal{L} = -\sum_{t=1}^{T-1} \log p(x_{t+1} \mid x_1, \ldots, x_t)$
            \item Used by: \textbf{TimesFM}, \textbf{Lag-Llama}, \textbf{TimeGPT}
        \end{itemize}
        \vspace{-1mm}
    \end{block}
    \vspace{-2.5mm}
    \begin{alertblock}{Paradigm 2: Masked Prediction (BERT-style)}
        \vspace{-1mm}
        \begin{itemize}\setlength{\itemsep}{0pt}
            \item Randomly mask 15--40\% of tokens, predict masked ones. Loss: $\mathcal{L} = -\sum_{t \in \mathcal{M}} \log p(x_t \mid x_{\setminus \mathcal{M}})$
            \item Used by: \textbf{Moirai}, \textbf{PatchTST} (self-supervised)
        \end{itemize}
        \vspace{-1mm}
    \end{alertblock}
    \vspace{-2.5mm}
    \begin{exampleblock}{Paradigm 3: Denoising / Seq2Seq (T5-style)}
        \vspace{-1mm}
        \begin{itemize}\setlength{\itemsep}{0pt}
            \item Encoder processes context, decoder generates forecast. Loss: cross-entropy over quantized tokens
            \item Used by: \textbf{Chronos}
        \end{itemize}
        \vspace{-1mm}
    \end{exampleblock}
    }
    \vspace{-3mm}
    \begin{center}
        \includegraphics[width=0.90\textwidth, height=0.13\textheight, keepaspectratio]{ch11_pretraining_objectives.pdf}
    \end{center}
    \end{cminipage}
\end{frame}

\begin{frame}{Data Augmentation for Pre-Training}
    \begin{cminipage}{0.95\textwidth}
    \vspace{-2mm}
    \begin{block}{KernelSynth (Ansari et al., 2024)}
        {\small
        \begin{itemize}\setlength{\itemsep}{0pt}
            \item Generate \textbf{synthetic} time series by composing Gaussian Process (GP) kernels:
            $$k(t, t') = k_{\text{RBF}}(t,t') + k_{\text{periodic}}(t,t') + k_{\text{linear}}(t,t') + \ldots$$
            \item Produces realistic patterns: trends, seasonality, changepoints, noise
            \item Chronos uses \textbf{billions} of synthetic series to augment real data
        \end{itemize}
        }
    \end{block}

    \vspace{-1mm}

    \begin{columns}[T]
        \column{0.48\textwidth}
        \begin{exampleblock}{TSMix (TimesFM)}
            {\scriptsize
            \begin{itemize}\setlength{\itemsep}{0pt}
                \item Mix existing series: $\tilde{x} = \lambda x_a + (1-\lambda) x_b$
                \item Augments diversity without needing GP expertise
                \item Combined with window slicing and frequency mixing
            \end{itemize}
            }
        \end{exampleblock}

        \column{0.48\textwidth}
        \begin{alertblock}{Why Augmentation Matters}
            {\scriptsize
            \begin{itemize}\setlength{\itemsep}{0pt}
                \item Prevents overfitting to specific domains
                \item Increases pattern diversity during pre-training
                \item Chronos: 50\% real data + 50\% synthetic $\Rightarrow$ \textbf{better} than 100\% real
            \end{itemize}
            }
        \end{alertblock}
    \end{columns}
    \end{cminipage}
\end{frame}

%=============================================================================
\section{Deep Dive into Models}
%=============================================================================

\begin{frame}{Chronos: Architecture}
    \begin{cminipage}{0.95\textwidth}
    \vspace{-3mm}
    {\scriptsize
    \begin{block}{Chronos (Ansari et al., 2024 --- Amazon)}
        \vspace{-1mm}
        \begin{itemize}\setlength{\itemsep}{0pt}
            \item \textbf{Architecture}: T5 encoder-decoder (pre-trained language model backbone)
            \item \textbf{Model sizes}: Mini (20M), Small (46M), Base (200M), Large (710M)
            \item \textbf{Tokenization}: Mean-scale normalization + quantization into \textbf{4096 bins}
            \item \textbf{Pre-training data}: 84B time points from 28 public datasets + KernelSynth
            \item \textbf{Output}: Categorical distribution over bins $\Rightarrow$ full probabilistic forecast
        \end{itemize}
        \vspace{-1mm}
    \end{block}
    \vspace{-2.5mm}
    \begin{columns}[T]
        \column{0.48\textwidth}
        \begin{exampleblock}{Key Innovations}
            \begin{itemize}\setlength{\itemsep}{0pt}
                \item Forecasting as \textbf{language modeling} over quantized tokens
                \item No architecture changes to T5 --- only tokenization is new
                \item Probabilistic: sample $N$ trajectories, compute quantiles
                \item Open-source, multiple sizes
            \end{itemize}
        \end{exampleblock}

        \column{0.48\textwidth}
        \begin{alertblock}{Limitations}
            \begin{itemize}\setlength{\itemsep}{0pt}
                \item Resolution limited by 4096 bins (quantization error)
                \item Univariate only (each variable separately)
                \item Context window: 512 tokens
                \item Slower inference than decoder-only (enc+dec)
            \end{itemize}
        \end{alertblock}
    \end{columns}
    }
    \quantlet{TSA\_ch11\_chronos}{https://github.com/QuantLet/TSA/tree/main/TSA_ch11/TSA_ch11_chronos}
    \end{cminipage}
\end{frame}

\begin{frame}{Chronos: How It Forecasts}
    \begin{cminipage}{0.95\textwidth}
    \vspace{-2mm}
    \begin{block}{Step-by-Step Pipeline}
        {\small
        \begin{enumerate}\setlength{\itemsep}{0pt}
            \item \textbf{Input}: raw time series $(x_1, x_2, \ldots, x_T)$
            \item \textbf{Normalize}: mean-scale $\tilde{x}_t = x_t / \bar{|x|}$
            \item \textbf{Quantize}: map each $\tilde{x}_t$ to one of 4096 bins via $\Phi$ (normal CDF)
            \item \textbf{Encode}: T5 encoder processes context tokens $\to$ hidden representations
            \item \textbf{Decode}: T5 decoder generates forecast tokens \textbf{autoregressively}
            \item \textbf{Sample}: repeat decoding $N$ times (e.g., $N=20$) with temperature sampling
            \item \textbf{Dequantize}: convert bin indices back to continuous values
            \item \textbf{Output}: median forecast + quantile intervals (10\%, 90\%)
        \end{enumerate}
        }
    \end{block}
    \vspace{-2mm}
    \begin{center}
        \includegraphics[width=0.95\textwidth, height=0.35\textheight, keepaspectratio]{ch11_chronos_pipeline.pdf}
    \end{center}
    \end{cminipage}
\end{frame}

\begin{frame}{Chronos: Benchmark Results}
    \begin{cminipage}{0.95\textwidth}
    \vspace{-2mm}
    \begin{block}{Performance on Monash Benchmark (Ansari et al., 2024)}
        {\scriptsize
        \begin{itemize}\setlength{\itemsep}{0pt}
            \item Evaluated on 27 datasets from the Monash Time Series Forecasting Archive
            \item Metrics: \textbf{WQL} (Weighted Quantile Loss) and \textbf{MASE} (Mean Absolute Scaled Error)
            \item Chronos-Large achieves best or second-best on 16/27 datasets \textbf{in zero-shot}
            \item Outperforms task-specific deep learning models (DeepAR, TFT, N-BEATS) on average
        \end{itemize}
        }
    \end{block}
    \vspace{-2mm}
    \begin{center}
        \includegraphics[width=0.95\textwidth, height=0.45\textheight, keepaspectratio]{ch11_chronos_benchmarks.pdf}
    \end{center}
    {\scriptsize WQL scores: lower is better. Chronos-Large competitive with or better than supervised baselines.}
    \end{cminipage}
\end{frame}

\begin{frame}{TimesFM: Architecture}
    \begin{cminipage}{0.95\textwidth}
    \vspace{-3mm}
    {\scriptsize
    \begin{block}{TimesFM (Das et al., 2024 --- Google Research)}
        \vspace{-1mm}
        \begin{itemize}\setlength{\itemsep}{0pt}
            \item \textbf{Architecture}: Decoder-only Transformer (200M parameters)
            \item \textbf{Input/output patching}: patch size $P = 32$, projects patch to model dimension
            \item \textbf{Pre-training data}: 100B real-world + synthetic data points
            \item \textbf{Multi-resolution}: trained on diverse frequencies (minutely to yearly)
            \item \textbf{Context window}: up to 512 patches ($= 16{,}384$ time steps)
        \end{itemize}
        \vspace{-1mm}
    \end{block}
    \vspace{-2.5mm}
    \begin{columns}[T]
        \column{0.48\textwidth}
        \begin{exampleblock}{Key Innovations}
            \begin{itemize}\setlength{\itemsep}{0pt}
                \item \textbf{Decoder-only}: simpler and faster than enc-dec
                \item \textbf{Output patches}: predict $P$ values at once
                \item \textbf{Frequency token}: encodes data frequency
                \item Point forecast + fine-tuning for quantiles
            \end{itemize}
        \end{exampleblock}

        \column{0.48\textwidth}
        \begin{alertblock}{Comparison with Chronos}
            \begin{itemize}\setlength{\itemsep}{0pt}
                \item Continuous values (no quantization)
                \item Faster inference (decoder-only, patch output)
                \item Less naturally probabilistic
                \item Closed weights until 2024, now open
            \end{itemize}
        \end{alertblock}
    \end{columns}
    }
    \quantlet{TSA\_ch11\_timesfm}{https://github.com/QuantLet/TSA/tree/main/TSA_ch11/TSA_ch11_timesfm}
    \end{cminipage}
\end{frame}

\begin{frame}{TimesFM: Zero-Shot Performance}
    \begin{cminipage}{0.95\textwidth}
    \vspace{-2mm}
    \begin{block}{Benchmark Results (Das et al., 2024)}
        {\small
        \begin{itemize}\setlength{\itemsep}{0pt}
            \item Evaluated on Monash, M3, M4, and long-horizon benchmarks (ETTh, Weather, Electricity)
            \item \textbf{Zero-shot} TimesFM beats supervised models on the majority of Monash datasets
            \item Particularly strong on \textbf{long-horizon} forecasting (96--720 steps)
            \item Competitive with task-specific PatchTST and iTransformer even without fine-tuning
        \end{itemize}
        }
    \end{block}
    \vspace{-2mm}
    \begin{center}
        \includegraphics[width=0.95\textwidth, height=0.43\textheight, keepaspectratio]{ch11_timesfm_benchmarks.pdf}
    \end{center}
    {\scriptsize TimesFM zero-shot MAE across benchmark datasets. Bars show relative improvement over best supervised baseline.}
    \end{cminipage}
\end{frame}

\begin{frame}{TimeGPT: API-First Foundation Model}
    \begin{cminipage}{0.95\textwidth}
    \vspace{-2mm}
    {\scriptsize
    \begin{block}{TimeGPT (Garza \& Mergenthaler-Canseco, 2024 --- Nixtla)}
        \begin{itemize}\setlength{\itemsep}{0pt}
            \item \textbf{Architecture}: Proprietary Transformer encoder-decoder
            \item \textbf{Pre-training}: largest known corpus --- \textbf{100+ billion} time points from diverse domains
            \item \textbf{Access}: API-only (no open weights), commercial product
            \item \textbf{Features}: zero-shot forecasting, conformal prediction intervals, anomaly detection, fine-tuning endpoint, exogenous variable support
        \end{itemize}
    \end{block}
    \vspace{-2mm}
    \begin{columns}[T]
        \column{0.48\textwidth}
        \begin{exampleblock}{Advantages}
            \begin{itemize}\setlength{\itemsep}{0pt}
                \item \textbf{Easiest to use}: 3 lines of Python code
                \item \textbf{Conformal prediction}: calibrated uncertainty
                \item \textbf{Exogenous variables}: unlike most FM
                \item Regular model updates by Nixtla team
            \end{itemize}
        \end{exampleblock}

        \column{0.48\textwidth}
        \begin{alertblock}{Limitations}
            \begin{itemize}\setlength{\itemsep}{0pt}
                \item \textbf{Black box}: proprietary, no reproducibility
                \item \textbf{API cost}: per-call pricing
                \item \textbf{Data privacy}: data sent to Nixtla servers
                \item \textbf{Vendor lock-in}: dependent on service availability
            \end{itemize}
        \end{alertblock}
    \end{columns}
    }
    \quantlet{TSA\_ch11\_timegpt}{https://github.com/QuantLet/TSA/tree/main/TSA_ch11/TSA_ch11_timegpt}
    \end{cminipage}
\end{frame}

\begin{frame}{Lag-Llama: LLM Backbone for Time Series}
    \begin{cminipage}{0.95\textwidth}
    \vspace{-2mm}
    {\scriptsize
    \begin{block}{Lag-Llama (Rasul et al., 2024)}
        \begin{itemize}\setlength{\itemsep}{0pt}
            \item \textbf{Architecture}: LLaMA decoder adapted for time series (352M parameters)
            \item \textbf{Key innovation}: \textbf{Lag features as token embeddings} --- lags $\{1, 2, \ldots, 7, 14, 28, \ldots\}$
            \item \textbf{Distribution head}: outputs Student-$t$ parameters $(\mu, \sigma, \nu)$ instead of token probabilities
            \item \textbf{Pre-training}: Multiple publicly available time series datasets
        \end{itemize}
    \end{block}
    \vspace{-2mm}
    \begin{columns}[T]
        \column{0.48\textwidth}
        \begin{exampleblock}{Why LLaMA for Time Series?}
            \begin{itemize}\setlength{\itemsep}{0pt}
                \item LLaMA backbone: proven at scale in NLP
                \item RoPE (Rotary Positional Encoding) for temporal positions
                \item GQA (Grouped Query Attention) for efficient inference
                \item Lag tokens capture known temporal structures (weekly, monthly)
            \end{itemize}
        \end{exampleblock}

        \column{0.48\textwidth}
        \begin{alertblock}{Distribution Head}
            \begin{itemize}\setlength{\itemsep}{0pt}
                \item Instead of softmax over vocabulary:
                $$p(x_{t+1} \mid \cdot) = \text{Student-}t(\mu, \sigma, \nu)$$
                \item Naturally handles heavy tails (financial data)
                \item Probabilistic: sample for prediction intervals
                \item Open-source and fine-tunable
            \end{itemize}
        \end{alertblock}
    \end{columns}
    }
    \quantlet{TSA\_ch11\_lagllama}{https://github.com/QuantLet/TSA/tree/main/TSA_ch11/TSA_ch11_lagllama}
    \end{cminipage}
\end{frame}

\begin{frame}{Moirai: Any-Variate, Any-Frequency}
    \begin{cminipage}{0.95\textwidth}
    \vspace{-2mm}
    {\scriptsize
    \begin{block}{Moirai (Woo et al., 2024 --- Salesforce)}
        \begin{itemize}\setlength{\itemsep}{0pt}
            \item \textbf{Architecture}: BERT-style masked encoder (14M, 91M, 311M parameters)
            \item \textbf{Key innovation}: handles \textbf{any number of variates}, \textbf{any frequency}, \textbf{any prediction length}
            \item \textbf{Pre-training}: 27 billion observations from \textbf{LOTSA} (Large-scale Open Time Series Archive, 9 domains)
            \item \textbf{Output}: mixture distribution (parameters of a flexible output distribution)
        \end{itemize}
    \end{block}
    \vspace{-2mm}
    \begin{columns}[T]
        \column{0.48\textwidth}
        \begin{exampleblock}{Universal Design}
            \begin{itemize}\setlength{\itemsep}{0pt}
                \item \textbf{Any-variate}: flattens multivariate into patches, uses variate IDs
                \item \textbf{Any-frequency}: frequency embeddings (hourly, daily, etc.)
                \item \textbf{Any-length}: masked attention for variable output lengths
                \item Closest to a ``universal'' TS foundation model
            \end{itemize}
        \end{exampleblock}

        \column{0.48\textwidth}
        \begin{alertblock}{LOTSA Dataset}
            \begin{itemize}\setlength{\itemsep}{0pt}
                \item 9 domains: energy, transport, finance, healthcare, nature, web, sales, climate, social
                \item 27B observations, curated for quality
                \item Open-source dataset collection
                \item Enables reproducible pre-training
            \end{itemize}
        \end{alertblock}
    \end{columns}
    }
    \quantlet{TSA\_ch11\_moirai}{https://github.com/QuantLet/TSA/tree/main/TSA_ch11/TSA_ch11_moirai}
    \end{cminipage}
\end{frame}

\begin{frame}{PatchTST: The Architectural Blueprint}
    \begin{cminipage}{0.95\textwidth}
    \vspace{-3mm}
    \begin{block}{PatchTST (Nie et al., 2023)}
        {\scriptsize
        \vspace{-1mm}
        \begin{itemize}\setlength{\itemsep}{0pt}
            \item \textbf{Architecture}: Transformer encoder with \textbf{channel-independent patching}
            \item Not a foundation model, but introduced two key ideas adopted by all later models:
            {\,}\textbf{(1)} Patching: group $P$ consecutive time points into one token;
            {\,}\textbf{(2)} Channel independence: each variate processed independently
            \item Parameters: 1M--10M (task-specific, much smaller than foundation models)
        \end{itemize}
        \vspace{-1mm}
        }
    \end{block}
    \vspace{-3mm}
    \begin{columns}[T]
        \column{0.50\textwidth}
        \begin{center}
            \includegraphics[width=0.95\textwidth, height=0.24\textheight, keepaspectratio]{ch11_patchtst_architecture.pdf}
        \end{center}

        \column{0.48\textwidth}
        \begin{exampleblock}{Legacy}
            {\scriptsize
            \begin{itemize}\setlength{\itemsep}{0pt}
                \item Showed Transformers can beat linear models on long-horizon forecasting
                \item Patching adopted by: TimesFM, Moirai, MOMENT
                \item Channel independence: Moirai, iTransformer
                \item Self-supervised variant: masked patching
            \end{itemize}
            }
        \end{exampleblock}
    \end{columns}
    \quantlet{TSA\_ch11\_patchtst}{https://github.com/QuantLet/TSA/tree/main/TSA_ch11/TSA_ch11_patchtst}
    \end{cminipage}
\end{frame}

\begin{frame}{Architecture Comparison Table}
    \begin{cminipage}{0.95\textwidth}
    \vspace{-3mm}
    {\scriptsize
    \begin{block}{Comprehensive Model Comparison}
    \vspace{-2mm}
    \begin{center}
    \begin{tabular}{l|l|r|l|l|l}
        \toprule
        \textbf{Model} & \textbf{Architecture} & \textbf{Params} & \textbf{Tokenization} & \textbf{Output} & \textbf{Pre-training} \\
        \midrule
        Chronos & T5 enc-dec & 20M--710M & Quantization (4096) & Categorical dist.\ & 84B pts + KernelSynth \\
        TimesFM & Decoder-only & 200M & Patching ($P\!=\!32$) & Point + quantiles & 100B real + synthetic \\
        TimeGPT & Enc-dec (prop.) & N/A & Proprietary & Point + conformal & 100B+ (proprietary) \\
        Lag-Llama & LLaMA decoder & 352M & Lag features & Student-$t$ & Public datasets \\
        Moirai & BERT encoder & 14M--311M & Patching + variate ID & Mixture dist.\ & 27B (LOTSA) \\
        PatchTST & Transformer enc & 1M--10M & Patching ($P\!=\!16$) & Point (MSE) & Task-specific \\
        \bottomrule
    \end{tabular}
    \end{center}
    \vspace{-2mm}
    \end{block}
    }
    \vspace{-2mm}
    \begin{exampleblock}{Key Takeaway}
        {\small No single architecture dominates. The choice depends on: (1)~whether you need probabilistic forecasts, (2)~multivariate support, (3)~compute budget, and (4)~whether open-source matters.}
    \end{exampleblock}
    \end{cminipage}
\end{frame}

\begin{frame}{The Scaling Hypothesis for Time Series}
    \begin{cminipage}{0.95\textwidth}
    \vspace{-3mm}
    \begin{block}{Does Scaling Work for Time Series?}
        {\scriptsize
        In NLP, \textbf{scaling laws} (Kaplan et al., 2020): loss $\propto N^{-\alpha}$ where $N$ = parameters. For time series: evidence is \textbf{mixed}.
        }
    \end{block}
    \vspace{-3mm}
    \begin{columns}[T]
        \column{0.45\textwidth}
        \begin{center}
            \includegraphics[width=0.95\textwidth, height=0.30\textheight, keepaspectratio]{ch11_scaling_laws.pdf}
        \end{center}

        \column{0.52\textwidth}
        {\scriptsize
        \begin{alertblock}{Why TS Scaling Differs from NLP}
            \begin{itemize}\setlength{\itemsep}{0pt}
                \item \textbf{Less data diversity}: text more diverse than TS
                \item \textbf{Simpler patterns}: many TS well-described by AR + seasonality
                \item \textbf{Diminishing returns}: Chronos-Large $\approx$ Chronos-Base
                \item \textbf{Data quality}: noisy TS limit learning
                \item \textbf{Evidence}: Chronos 710M $\approx$ 200M; TimesFM 200M competitive with larger models; data quality $>$ parameter count
            \end{itemize}
        \end{alertblock}
        }
    \end{columns}
    \end{cminipage}
\end{frame}

\begin{frame}{When Foundation Models Fail}
    \begin{cminipage}{0.95\textwidth}
    \vspace{-3mm}
    \begin{alertblock}{Failure Modes to Watch For}
        {\scriptsize
        \begin{enumerate}\setlength{\itemsep}{0pt}
            \item \textbf{Distribution shift}: Financial regime changes (COVID crash, rate hikes). Pre-training data may not cover such regimes.
            \item \textbf{Missing exogenous variables}: Weather affects energy demand; policy changes affect economics. Most foundation models are \textbf{univariate}.
            \item \textbf{Rare frequencies}: Irregular sampling, event-driven or tick-by-tick data. Pre-training favors regular frequencies.
            \item \textbf{Very short series}: Fewer than $\sim$50 observations $\Rightarrow$ insufficient context. Classical methods (ETS, ARIMA) often dominate.
            \item \textbf{Multivariate dependencies}: Cross-variable correlations (e.g., VAR) ignored by univariate models. Only Moirai attempts multivariate support.
        \end{enumerate}
        }
    \end{alertblock}
    \vspace{-3mm}
    \begin{exampleblock}{Rule of Thumb}
        {\scriptsize If your data has \textbf{known causal structure}, \textbf{domain-specific features}, or \textbf{regime changes}, consider combining foundation model forecasts with domain-informed models (ensemble or hierarchical).}
    \end{exampleblock}
    \end{cminipage}
\end{frame}

%=============================================================================
\section{Zero-Shot vs Fine-Tuning}
%=============================================================================

\begin{frame}{Zero-Shot: When and Why}
    \begin{cminipage}{0.95\textwidth}
    \vspace{-2mm}
    \begin{block}{Zero-Shot Forecasting}
        {\small Use the pre-trained model \textbf{directly} on a new series without any training.}
    \end{block}
    \vspace{-2mm}
    \begin{columns}[T]
        \column{0.48\textwidth}
        \begin{exampleblock}{When to Use Zero-Shot}
            {\scriptsize
            \begin{itemize}\setlength{\itemsep}{0pt}
                \item \textbf{Limited labeled data}: new product launch, cold-start
                \item \textbf{Rapid prototyping}: quick baseline in minutes
                \item \textbf{Cross-domain deployment}: same model for finance, retail, energy
                \item \textbf{Many series}: 10{,}000+ series to forecast, no time to train each
                \item \textbf{No ML expertise}: simple API call
            \end{itemize}
            }
        \end{exampleblock}

        \column{0.48\textwidth}
        \begin{center}
            \includegraphics[width=0.95\textwidth, height=0.45\textheight, keepaspectratio]{ch11_zeroshot_vs_finetuning.pdf}
        \end{center}
    \end{columns}
    \vspace{-2mm}
    \begin{alertblock}{Surprising Finding}
        {\small Ansari et al.\ (2024): Chronos zero-shot outperforms task-specific models (DeepAR, N-BEATS) that were \textit{trained on the target data}. Pre-training transfers temporal knowledge effectively.}
    \end{alertblock}
    \end{cminipage}
\end{frame}

\begin{frame}{Fine-Tuning: Domain Adaptation}
    \begin{cminipage}{0.95\textwidth}
    \vspace{-2mm}
    \begin{block}{Full Fine-Tuning vs LoRA}
        {\small
        \begin{itemize}\setlength{\itemsep}{0pt}
            \item \textbf{Full fine-tuning}: update all parameters on domain data (expensive, risk of forgetting)
            \item \textbf{LoRA} (Low-Rank Adaptation, Hu et al., 2022): freeze base weights, add small trainable matrices
            $$W' = W + \Delta W = W + BA, \quad B \in \mathbb{R}^{d \times r},\; A \in \mathbb{R}^{r \times d},\; r \ll d$$
            \item LoRA: 0.1--1\% of trainable parameters, similar performance, much cheaper
        \end{itemize}
        }
    \end{block}
    \vspace{-2mm}
    \begin{columns}[T]
        \column{0.48\textwidth}
        \begin{alertblock}{Risks of Fine-Tuning}
            {\scriptsize
            \begin{itemize}\setlength{\itemsep}{0pt}
                \item \textbf{Catastrophic forgetting}: model loses pre-trained knowledge
                \item \textbf{Overfitting}: small domain dataset $\Rightarrow$ poor generalization
                \item \textbf{Hyperparameter sensitivity}: learning rate, epochs, warmup critical
            \end{itemize}
            }
        \end{alertblock}

        \column{0.48\textwidth}
        \begin{exampleblock}{Best Practices}
            {\scriptsize
            \begin{itemize}\setlength{\itemsep}{0pt}
                \item Start with \textbf{low learning rate} ($10^{-5}$--$10^{-4}$)
                \item Use \textbf{learning rate warmup} (500--1000 steps)
                \item \textbf{Early stopping} on validation loss
                \item \textbf{Freeze early layers}, fine-tune later layers
                \item Prefer \textbf{LoRA} for small datasets ($<$10K series)
            \end{itemize}
            }
        \end{exampleblock}
    \end{columns}
    \end{cminipage}
\end{frame}

\begin{frame}{Few-Shot and In-Context Learning}
    \begin{cminipage}{0.95\textwidth}
    \vspace{-2mm}
    \begin{block}{Context Length as ``Few-Shot'' for Time Series}
        {\small
        \begin{itemize}\setlength{\itemsep}{0pt}
            \item In NLP, ``few-shot'' = provide examples in the prompt. In TS, ``context'' = historical observations
            \item \textbf{Longer context $\to$ better forecasts} (up to the model's context window)
            \item The model conditions on historical patterns to infer trend, seasonality, and noise structure
        \end{itemize}
        }
    \end{block}
    \vspace{-2mm}
    \begin{columns}[T]
        \column{0.50\textwidth}
        \begin{center}
            \includegraphics[width=0.95\textwidth, height=0.38\textheight, keepaspectratio]{ch11_context_length_effect.pdf}
        \end{center}
        {\scriptsize Forecast error vs.\ context length: diminishing returns beyond the model's window.}

        \column{0.48\textwidth}
        {\scriptsize
        \begin{exampleblock}{Context Windows by Model}
            \begin{itemize}\setlength{\itemsep}{0pt}
                \item Chronos: 512 tokens ($\approx$ 512 time steps)
                \item TimesFM: 512 patches ($\approx$ 16{,}384 steps)
                \item Moirai: 5000 tokens (configurable)
                \item Lag-Llama: 1024 steps (with lags extending further)
                \item TimeGPT: configurable via API
            \end{itemize}
        \end{exampleblock}
        \vspace{-2mm}
        \begin{alertblock}{Practical Tip}
            Provide as much context as the model allows. For daily data, $\geq$1 year is ideal. For monthly data, $\geq$5 years captures seasonal cycles.
        \end{alertblock}
        }
    \end{columns}
    \end{cminipage}
\end{frame}

\begin{frame}{Decision Framework}
    \begin{cminipage}{0.95\textwidth}
    \vspace{-2mm}
    \begin{block}{When to Choose Zero-Shot, Fine-Tuning, or Classical}
        {\small
        \begin{itemize}\setlength{\itemsep}{0pt}
            \item \textbf{Step 1}: Do you have labeled domain data? No $\to$ Zero-shot
            \item \textbf{Step 2}: Is zero-shot performance sufficient? Yes $\to$ Deploy; No $\to$ Fine-tune
            \item \textbf{Step 3}: Is the domain very different from pre-training data? Yes $\to$ Full fine-tune; No $\to$ LoRA
            \item \textbf{Step 4}: Need interpretability or causal reasoning? Yes $\to$ Classical models (ARIMA, VAR)
        \end{itemize}
        }
    \end{block}
    \vspace{-2mm}
    \begin{columns}[T]
        \column{0.50\textwidth}
        \begin{center}
            \includegraphics[width=0.95\textwidth, height=0.38\textheight, keepaspectratio]{ch11_decision_flowchart.pdf}
        \end{center}

        \column{0.48\textwidth}
        \begin{exampleblock}{Rules of Thumb}
            {\scriptsize
            \begin{itemize}\setlength{\itemsep}{0pt}
                \item $<$50 observations: ARIMA / ETS
                \item 50--500 observations: zero-shot foundation model
                \item 500--10K observations: fine-tuned foundation model
                \item $>$10K observations: task-specific deep learning may match or beat FM
                \item Always benchmark against ARIMA as baseline
            \end{itemize}
            }
        \end{exampleblock}
    \end{columns}
    \end{cminipage}
\end{frame}

%=============================================================================
\section{Practical Use Cases}
%=============================================================================

\begin{frame}{Use Case 1: EUR/RON with All Foundation Models}
    \begin{cminipage}{0.95\textwidth}
    \vspace{-2mm}
    \begin{block}{Extending the Chapter 8 Case Study}
        {\small
        \begin{itemize}\setlength{\itemsep}{0pt}
            \item \textbf{Data}: EUR/RON daily exchange rate (2020--2024, $\sim$1{,}200 observations)
            \item \textbf{Task}: 30-day-ahead forecast, 80/20 temporal split
            \item \textbf{Models compared}: ARIMA, ARFIMA, Random Forest, LSTM, \textbf{Chronos}, \textbf{TimesFM}, \textbf{Lag-Llama}, \textbf{Moirai}, \textbf{TimeGPT}
            \item Foundation models used in \textbf{zero-shot} mode (no training on EUR/RON data)
        \end{itemize}
        }
    \end{block}
    \vspace{-2mm}
    \begin{center}
        \includegraphics[width=0.95\textwidth, height=0.43\textheight, keepaspectratio]{ch11_eurron_all_models.pdf}
    \end{center}
    \quantlet{TSA\_ch11\_eurron\_foundation}{https://github.com/QuantLet/TSA/tree/main/TSA_ch11/TSA_ch11_eurron_foundation}
    \end{cminipage}
\end{frame}

\begin{frame}{EUR/RON: Detailed Results}
    \begin{cminipage}{0.95\textwidth}
    \vspace{-3mm}
    {\scriptsize
    \begin{block}{Forecast Performance Comparison}
    \vspace{-2mm}
    \begin{center}
    \begin{tabular}{l|c|c|c|c}
        \toprule
        \textbf{Model} & \textbf{RMSE} & \textbf{MAE} & \textbf{MAPE (\%)} & \textbf{Training Time} \\
        \midrule
        ARIMA(1,1,1) & --- & --- & --- & $<$1s \\
        ARFIMA & --- & --- & --- & $<$1s \\
        Random Forest & --- & --- & --- & $\sim$2s \\
        LSTM (2-layer) & --- & --- & --- & $\sim$60s \\
        Chronos-Base & --- & --- & --- & 0s (zero-shot) \\
        TimesFM & --- & --- & --- & 0s (zero-shot) \\
        Lag-Llama & --- & --- & --- & 0s (zero-shot) \\
        Moirai-Base & --- & --- & --- & 0s (zero-shot) \\
        TimeGPT & --- & --- & --- & 0s (API call) \\
        \bottomrule
    \end{tabular}
    \end{center}
    \vspace{-2mm}
    {\scriptsize Note: Metric values (---) populated at runtime from the Quantlet code. See the companion notebook for exact numbers.}
    \end{block}
    }
    \vspace{-2mm}
    \begin{center}
        \includegraphics[width=0.95\textwidth, height=0.32\textheight, keepaspectratio]{ch11_eurron_predictions_detail.pdf}
    \end{center}
    \end{cminipage}
\end{frame}

\begin{frame}{Use Case 2: Energy Demand Forecasting}
    \begin{cminipage}{0.95\textwidth}
    \vspace{-2mm}
    \begin{block}{German Electricity Demand (Hourly, Multi-Seasonal)}
        {\small
        \begin{itemize}\setlength{\itemsep}{0pt}
            \item \textbf{Data}: German hourly electricity consumption (2015--2023, 70K+ observations)
            \item \textbf{Challenge}: multiple seasonalities (24h daily + 168h weekly + annual)
            \item \textbf{Models}: TBATS (Ch9), Prophet (Ch9), LSTM (Ch8), Chronos, TimesFM, Moirai
            \item Foundation models handle multi-seasonal patterns \textbf{without explicit specification}
        \end{itemize}
        }
    \end{block}
    \vspace{-2mm}
    \begin{center}
        \includegraphics[width=0.95\textwidth, height=0.43\textheight, keepaspectratio]{ch11_energy_foundation.pdf}
    \end{center}
    \quantlet{TSA\_ch11\_energy\_foundation}{https://github.com/QuantLet/TSA/tree/main/TSA_ch11/TSA_ch11_energy_foundation}
    \end{cminipage}
\end{frame}

\begin{frame}{Use Case 3: Financial Volatility}
    \begin{cminipage}{0.95\textwidth}
    \vspace{-2mm}
    \begin{block}{S\&P 500 Realized Volatility}
        {\small
        \begin{itemize}\setlength{\itemsep}{0pt}
            \item \textbf{Data}: S\&P 500 daily realized volatility (2000--2024)
            \item \textbf{Challenge}: Clustering, heavy tails, regime changes (2008 crisis, 2020 COVID)
            \item \textbf{Models}: GARCH(1,1) (Ch5), HAR-RV, Chronos, TimesFM, Lag-Llama
            \item \textbf{Question}: Can foundation models capture volatility dynamics without GARCH structure?
        \end{itemize}
        }
    \end{block}
    \vspace{-2mm}
    \begin{center}
        \includegraphics[width=0.95\textwidth, height=0.43\textheight, keepaspectratio]{ch11_volatility_comparison.pdf}
    \end{center}
    \quantlet{TSA\_ch11\_volatility\_foundation}{https://github.com/QuantLet/TSA/tree/main/TSA_ch11/TSA_ch11_volatility_foundation}
    \end{cminipage}
\end{frame}

\begin{frame}[fragile]{Python Code Patterns}
    \begin{cminipage}{0.95\textwidth}
    \vspace{-3mm}
    {\scriptsize
    \begin{columns}[T]
        \column{0.48\textwidth}
        \begin{block}{Chronos (3 lines)}
            \begin{verbatim}
from chronos import ChronosPipeline
pipe = ChronosPipeline.from_pretrained(
    "amazon/chronos-t5-base")
forecast = pipe.predict(
    context=torch.tensor(data),
    prediction_length=30)
            \end{verbatim}
        \end{block}
        \vspace{-3mm}
        \begin{block}{TimesFM (3 lines)}
            \begin{verbatim}
import timesfm
tfm = timesfm.TimesFm(hparams=
    timesfm.TimesFmHparams())
tfm.load_from_checkpoint()
preds = tfm.forecast(data)
            \end{verbatim}
        \end{block}

        \column{0.48\textwidth}
        \begin{block}{TimeGPT (3 lines)}
            \begin{verbatim}
from nixtla import NixtlaClient
nixtla = NixtlaClient(api_key="...")
fcst = nixtla.forecast(
    df=df, h=30, model="timegpt-1")
            \end{verbatim}
        \end{block}
        \vspace{-3mm}
        \begin{block}{Lag-Llama (4 lines)}
            \begin{verbatim}
from lag_llama.gluon.estimator import
    LagLlamaEstimator
estimator = LagLlamaEstimator(
    prediction_length=30)
predictor = estimator.train(train_ds)
forecasts = list(
    predictor.predict(test_ds))
            \end{verbatim}
        \end{block}
    \end{columns}
    }
    \end{cminipage}
\end{frame}

%=============================================================================
\section{Benchmarks and Model Selection}
%=============================================================================

\begin{frame}{Major Benchmarks}
    \begin{cminipage}{0.95\textwidth}
    \vspace{-2mm}
    {\scriptsize
    \begin{block}{Monash Time Series Forecasting Archive}
        \begin{itemize}\setlength{\itemsep}{0pt}
            \item \textbf{30K+ series} across 30 datasets (tourism, electricity, traffic, weather, finance, etc.)
            \item Standard evaluation protocol: fixed train/test splits, WQL and MASE metrics
            \item Most widely used benchmark for foundation models (Chronos, TimesFM, Moirai all report here)
        \end{itemize}
    \end{block}
    \vspace{-2mm}
    \begin{columns}[T]
        \column{0.48\textwidth}
        \begin{exampleblock}{M-Competitions}
            \begin{itemize}\setlength{\itemsep}{0pt}
                \item \textbf{M3} (2000): 3{,}003 series, monthly/quarterly/yearly
                \item \textbf{M4} (2018): 100{,}000 series, 6 frequencies
                \item \textbf{M5} (2020): Walmart sales, hierarchical
                \item Finding: simpler methods often competitive
            \end{itemize}
        \end{exampleblock}

        \column{0.48\textwidth}
        \begin{alertblock}{GiftEval (2024)}
            \begin{itemize}\setlength{\itemsep}{0pt}
                \item \textbf{General Time Series Forecasting Evaluation}
                \item Designed specifically for foundation model evaluation
                \item Diverse frequencies, domains, horizons
                \item Calibration and reliability metrics
            \end{itemize}
        \end{alertblock}
    \end{columns}
    }
    \end{cminipage}
\end{frame}

\begin{frame}{Aggregated Results Across Benchmarks}
    \begin{cminipage}{0.95\textwidth}
    \vspace{-2mm}
    \begin{block}{Cross-Benchmark Performance Summary}
        {\small
        \begin{itemize}\setlength{\itemsep}{0pt}
            \item Aggregated rankings across Monash, M3, M4, ETT, Weather, Electricity, and Traffic
            \item No single model wins everywhere --- \textbf{model selection depends on the use case}
            \item Foundation models are \textbf{most impressive in zero-shot}, often matching supervised baselines
        \end{itemize}
        }
    \end{block}
    \vspace{-2mm}
    \begin{center}
        \includegraphics[width=0.95\textwidth, height=0.43\textheight, keepaspectratio]{ch11_aggregated_benchmarks.pdf}
    \end{center}
    {\scriptsize Heatmap: model (rows) $\times$ benchmark (columns). Color = relative ranking (green = top, red = bottom).}
    \quantlet{TSA\_ch11\_benchmarks}{https://github.com/QuantLet/TSA/tree/main/TSA_ch11/TSA_ch11_benchmarks}
    \end{cminipage}
\end{frame}

\begin{frame}{When to Use What: Decision Table}
    \begin{cminipage}{0.95\textwidth}
    \vspace{-3mm}
    {\scriptsize
    \begin{block}{Scenario-Based Recommendations}
    \vspace{-2mm}
    \begin{center}
    \begin{tabular}{l|c|c|c|c|c}
        \toprule
        \textbf{Scenario} & \textbf{ARIMA} & \textbf{Chronos} & \textbf{TimesFM} & \textbf{Moirai} & \textbf{TimeGPT} \\
        \midrule
        Limited data ($<$100 obs) & $\checkmark\checkmark$ & $\checkmark$ & $\checkmark$ & $\sim$ & $\checkmark$ \\
        Real-time / low latency & $\checkmark\checkmark$ & $\sim$ & $\checkmark$ & $\sim$ & $\checkmark$ \\
        Interpretability needed & $\checkmark\checkmark$ & $\times$ & $\times$ & $\times$ & $\times$ \\
        Multivariate dependencies & VAR & $\times$ & $\times$ & $\checkmark\checkmark$ & $\sim$ \\
        Many series (10K+) & $\sim$ & $\checkmark\checkmark$ & $\checkmark\checkmark$ & $\checkmark\checkmark$ & $\checkmark$ \\
        No ML expertise & $\checkmark$ & $\checkmark$ & $\sim$ & $\sim$ & $\checkmark\checkmark$ \\
        Open source required & $\checkmark\checkmark$ & $\checkmark\checkmark$ & $\checkmark$ & $\checkmark\checkmark$ & $\times$ \\
        Probabilistic forecasts & $\checkmark$ & $\checkmark\checkmark$ & $\sim$ & $\checkmark\checkmark$ & $\checkmark$ \\
        \bottomrule
    \end{tabular}
    \end{center}
    \vspace{-2mm}
    {\scriptsize $\checkmark\checkmark$ = excellent, $\checkmark$ = good, $\sim$ = acceptable, $\times$ = poor fit}
    \end{block}
    }
    \end{cminipage}
\end{frame}

\begin{frame}{Computational Cost Analysis}
    \begin{cminipage}{0.95\textwidth}
    \vspace{-2mm}
    \begin{block}{Pareto Frontier: Accuracy vs.\ Cost}
        {\small
        \begin{itemize}\setlength{\itemsep}{0pt}
            \item Foundation models trade \textbf{compute for convenience}: no training needed, but inference is heavier
            \item The ``best'' model lies on the Pareto frontier of accuracy and cost
        \end{itemize}
        }
    \end{block}
    \vspace{-2mm}
    \begin{columns}[T]
        \column{0.50\textwidth}
        \begin{center}
            \includegraphics[width=0.95\textwidth, height=0.40\textheight, keepaspectratio]{ch11_accuracy_vs_cost.pdf}
        \end{center}

        \column{0.48\textwidth}
        {\scriptsize
        \begin{alertblock}{Cost Considerations}
            \begin{itemize}\setlength{\itemsep}{0pt}
                \item \textbf{ARIMA}: CPU, milliseconds, free
                \item \textbf{Chronos-Mini}: 1 GPU, $\sim$0.1s/series
                \item \textbf{Chronos-Large}: 1 GPU (A100), $\sim$1s/series
                \item \textbf{TimesFM}: 1 GPU, $\sim$0.05s/series (fast)
                \item \textbf{TimeGPT}: API, \$0.01--\$0.10/call
                \item \textbf{Pre-training from scratch}: 100+ GPU-hours (Chronos-Large)
            \end{itemize}
        \end{alertblock}
        }
    \end{columns}
    \end{cminipage}
\end{frame}

%=============================================================================
\section{Limitations, Ethics, and Future}
%=============================================================================

\begin{frame}{Current Limitations}
    \begin{cminipage}{0.95\textwidth}
    \vspace{-2mm}
    \begin{alertblock}{Five Key Limitations of Foundation Models for Time Series}
        {\small
        \begin{enumerate}\setlength{\itemsep}{2pt}
            \item \textbf{Most models are univariate only}: Cannot capture cross-variable interactions (e.g., price-volume, temperature-demand). Only Moirai attempts multivariate.
            \item \textbf{Limited context windows (512--4096 tokens)}: Cannot leverage very long histories. For daily data with yearly seasonality, 512 steps $\approx$ 1.5 years only.
            \item \textbf{No causal reasoning}: Models learn correlations, not causation. Cannot answer ``what if interest rates rise by 1\%?'' --- no structural model underneath.
            \item \textbf{Black box nature}: Attention weights are not easily interpretable. Difficult to explain \textit{why} a particular forecast was made --- problematic for regulated industries.
            \item \textbf{``Hallucination'' in probabilistic forecasts}: Generated trajectories may exhibit unrealistic patterns (e.g., negative prices, explosive growth) that violate domain constraints.
        \end{enumerate}
        }
    \end{alertblock}
    \vspace{-2mm}
    \begin{exampleblock}{Mitigation}
        {\small Post-processing (clipping, reconciliation), ensemble with classical models, conformal calibration.}
    \end{exampleblock}
    \end{cminipage}
\end{frame}

\begin{frame}{Ethical Considerations}
    \begin{cminipage}{0.95\textwidth}
    \vspace{-2mm}
    \begin{block}{Data Provenance}
        {\small
        \begin{itemize}\setlength{\itemsep}{0pt}
            \item Pre-training data may include \textbf{private or proprietary} time series (e.g., internal sales data scraped from APIs)
            \item Lack of transparency: most models do not fully disclose training data composition
            \item LOTSA (Moirai) is a positive example: open, curated, documented sources
        \end{itemize}
        }
    \end{block}
    \vspace{-2mm}
    \begin{columns}[T]
        \column{0.48\textwidth}
        \begin{alertblock}{Fairness and Bias}
            {\scriptsize
            \begin{itemize}\setlength{\itemsep}{0pt}
                \item Pre-training skews toward \textbf{data-rich domains} (energy, retail)
                \item Under-represented: healthcare in developing countries, agriculture in Africa
                \item Biased training $\to$ biased forecasts $\to$ biased decisions
            \end{itemize}
            }
        \end{alertblock}

        \column{0.48\textwidth}
        \begin{alertblock}{Environmental and Regulatory}
            {\scriptsize
            \begin{itemize}\setlength{\itemsep}{0pt}
                \item \textbf{Carbon footprint}: pre-training Chronos-Large $\approx$ 100+ GPU-hours on A100s
                \item \textbf{EU AI Act}: high-risk AI systems require explainability --- foundation models may not qualify
                \item \textbf{Accountability}: who is responsible when a forecast-based decision fails?
            \end{itemize}
            }
        \end{alertblock}
    \end{columns}
    \end{cminipage}
\end{frame}

\begin{frame}{Future Directions}
    \begin{cminipage}{0.95\textwidth}
    \vspace{-2mm}
    \begin{block}{Where Is the Field Heading?}
        {\small
        \begin{enumerate}\setlength{\itemsep}{2pt}
            \item \textbf{Multimodal TS models}: Combine time series with text (news, reports), images (satellite), and metadata. Example: forecast sales using both historical data and product descriptions.
            \item \textbf{Agentic forecasting}: LLM agents that \textit{select, run, and evaluate} forecasting models autonomously. The LLM orchestrates the pipeline rather than producing the forecast directly.
            \item \textbf{Continual pre-training}: Update foundation models with new data without full retraining. Critical for adapting to distribution shifts (post-COVID, post-pandemic).
            \item \textbf{Multivariate foundation models}: Move beyond univariate. Learn cross-variable dynamics, hierarchical structures, and spatial correlations.
            \item \textbf{Domain-specific foundation models}: FinanceFM, EnergyFM, HealthFM --- smaller, specialized, more accurate within their domain.
        \end{enumerate}
        }
    \end{block}
    \vspace{-2mm}
    \begin{exampleblock}{The Next 3--5 Years}
        {\small Foundation models will likely become the \textbf{default starting point} for forecasting, much as pre-trained word embeddings became standard in NLP. Classical methods will remain essential for interpretability, causal inference, and small-data regimes.}
    \end{exampleblock}
    \end{cminipage}
\end{frame}

\begin{frame}{A Balanced Perspective}
    \begin{cminipage}{0.95\textwidth}
    \vspace{-2mm}
    \begin{alertblock}{When to Use Classical Models}
        {\small
        \begin{itemize}\setlength{\itemsep}{0pt}
            \item \textbf{Interpretability required}: regulatory settings, scientific research, policy analysis
            \item \textbf{Causal inference}: structural models (VAR, SVAR) for ``what if'' analysis
            \item \textbf{Small data}: fewer than $\sim$50 observations, seasonal structure known a priori
            \item \textbf{Computational constraints}: edge devices, real-time systems, no GPU access
            \item \textbf{Explainability mandate}: EU AI Act, financial regulation (MiFID II, Basel III)
        \end{itemize}
        }
    \end{alertblock}
    \vspace{-2mm}
    \begin{exampleblock}{When to Use Foundation Models}
        {\small
        \begin{itemize}\setlength{\itemsep}{0pt}
            \item \textbf{Rapid prototyping}: get a strong baseline in minutes, not days
            \item \textbf{Cross-domain}: same model for retail, energy, finance --- no domain expertise needed
            \item \textbf{Zero-shot / cold-start}: new product, new market, no historical data
            \item \textbf{Scale}: 10{,}000+ series to forecast --- foundation models amortize the cost
            \item \textbf{Probabilistic forecasts}: native uncertainty quantification (quantiles, prediction intervals)
        \end{itemize}
        }
    \end{exampleblock}
    \vspace{-2mm}
    \begin{block}{}
        {\small \textbf{Bottom line}: Classical and foundation models are \textbf{complementary}, not competing. The best practitioners use both.}
    \end{block}
    \end{cminipage}
\end{frame}

%=============================================================================
\section{AI Use Case}
%=============================================================================

\begin{frame}{AI Exercise: Foundation Model Selection}
    \begin{cminipage}{0.95\textwidth}
    \vspace{-3mm}
    \begin{block}{\footnotesize Prompt to test in ChatGPT / Claude / Copilot}
        {\footnotesize
        ``Compare ARIMA, Prophet, and a foundation model (Chronos) for forecasting monthly airline passengers (the classic Box-Jenkins dataset, 144 observations). For each model: (1)~fit the model, (2)~forecast 24 months ahead, (3)~plot actual vs predicted with 95\% confidence intervals, (4)~compute RMSE and MAPE. Which model would you recommend and why? Provide complete Python code.''
        }
    \end{block}
    \vspace{-2mm}
    {\footnotesize
    \textbf{Exercise}:
    \begin{enumerate}\setlength{\itemsep}{0pt}
        \item Run the prompt in an LLM of your choice and critically analyze the response.
        \item Does the code properly handle the \textbf{multiplicative seasonality} in airline data?
        \item How does Chronos handle seasonality \textbf{without being told} the period is 12?
        \item Are confidence intervals from Chronos (quantile-based) comparable to ARIMA's (parametric)?
        \item Does the LLM discuss the \textbf{trade-off}: ARIMA is interpretable, Chronos is zero-shot?
    \end{enumerate}
    }
    \vspace{-2mm}
    \begin{alertblock}{}
        {\footnotesize \textbf{Warning}: AI-generated code may run without errors and look professional. \textit{That does not mean it is correct.} Always validate against known benchmarks.}
    \end{alertblock}
    \end{cminipage}
\end{frame}

%=============================================================================
\section{Summary}
%=============================================================================

\begin{frame}{Key Takeaways}
    \begin{cminipage}{0.95\textwidth}
    \begin{block}{What We Learned}
        \begin{itemize}\setlength{\itemsep}{2pt}
            \item \textbf{Transformers} replace recurrence with self-attention: parallel, scalable, long-range dependencies
            \item \textbf{Foundation models} are pre-trained on billions of diverse time points and deployed zero-shot
            \item \textbf{Tokenization} strategies (patching, quantization) bridge continuous TS and discrete Transformer inputs
            \item \textbf{Chronos} (T5, quantized), \textbf{TimesFM} (decoder, patched), \textbf{Lag-Llama} (LLaMA, lags), \textbf{Moirai} (BERT, any-variate) each have distinct design philosophies
            \item \textbf{Zero-shot} performance is surprisingly strong --- often matching supervised baselines
            \item \textbf{Fine-tuning} (especially LoRA) helps when domain data is available but must be done carefully
        \end{itemize}
    \end{block}

    \vspace{0.1cm}

    \begin{alertblock}{Practical Recommendations}
        \begin{itemize}\setlength{\itemsep}{0pt}
            \item Always benchmark against \textbf{ARIMA/ETS} as a classical baseline
            \item Use \textbf{foundation models} for rapid prototyping and when domain expertise is limited
            \item Classical and foundation models are \textbf{complementary} --- combine them via ensembles
        \end{itemize}
    \end{alertblock}
    \end{cminipage}
\end{frame}

%=============================================================================
\section{Quiz}
%=============================================================================

\begin{frame}{Question 1}
    \begin{cminipage}{0.95\textwidth}
    \begin{alertblock}{Question}
        \begin{itemize}\setlength{\itemsep}{0pt}
            \item What is the key advantage of self-attention over recurrence in sequence modeling?
        \end{itemize}
    \end{alertblock}

    \vspace{0.3cm}

    \begin{block}{Answer Choices}

        \textcolor{MainBlue}{\textbf{(A)}} Self-attention uses fewer parameters than RNNs\\[3pt]

        \textcolor{MainBlue}{\textbf{(B)}} Self-attention processes all positions in parallel with $O(1)$ path length between any two positions\\[3pt]

        \textcolor{MainBlue}{\textbf{(C)}} Self-attention does not require positional encoding\\[3pt]

        \textcolor{MainBlue}{\textbf{(D)}} Self-attention is always faster than recurrence regardless of sequence length

    \end{block}
    \end{cminipage}
\end{frame}

\begin{frame}{Question 1: Answer}
    \begin{cminipage}{0.95\textwidth}
    \vspace{-0.2cm}
    \begin{center}
        \includegraphics[width=0.98\textwidth, height=0.58\textheight, keepaspectratio]{ch11_quiz1_attention_vs_rnn.pdf}
    \end{center}
    \vspace{-3mm}
    {\small
    \begin{exampleblock}{Answer: (B)}
    \begin{itemize}\setlength{\itemsep}{0pt}
        \item Self-attention connects every position to every other position directly (path length $O(1)$), enabling full parallelization. RNNs require $O(T)$ sequential steps, causing vanishing gradients and slow training.
    \end{itemize}
    \end{exampleblock}
    }
    \hfill\quantlet{TSA\_ch11\_quiz1\_attention\_vs\_rnn}{https://github.com/QuantLet/TSA/tree/main/TSA_ch11/TSA_ch11_quiz1_attention_vs_rnn}
    \end{cminipage}
\end{frame}

\begin{frame}{Question 2}
    \begin{cminipage}{0.95\textwidth}
    \begin{alertblock}{Question}
        \begin{itemize}\setlength{\itemsep}{0pt}
            \item Which tokenization strategy does Chronos use to convert continuous time series into discrete tokens?
        \end{itemize}
    \end{alertblock}

    \vspace{0.3cm}

    \begin{block}{Answer Choices}

        \textcolor{MainBlue}{\textbf{(A)}} Patching: grouping consecutive time points into patches\\[3pt]

        \textcolor{MainBlue}{\textbf{(B)}} Scaling and quantization into discrete bins using the normal CDF\\[3pt]

        \textcolor{MainBlue}{\textbf{(C)}} Direct linear embedding of raw values\\[3pt]

        \textcolor{MainBlue}{\textbf{(D)}} Byte-pair encoding (BPE) borrowed from NLP

    \end{block}
    \end{cminipage}
\end{frame}

\begin{frame}{Question 2: Answer}
    \begin{cminipage}{0.95\textwidth}
    \vspace{-0.2cm}
    \begin{center}
        \includegraphics[width=0.98\textwidth, height=0.58\textheight, keepaspectratio]{ch11_quiz2_tokenization.pdf}
    \end{center}
    \vspace{-3mm}
    {\small
    \begin{exampleblock}{Answer: (B)}
    \begin{itemize}\setlength{\itemsep}{0pt}
        \item Chronos normalizes values (mean-scaling), then maps to 4096 bins via $b_t = \lfloor (B-1) \cdot \Phi(\tilde{x}_t) \rfloor$. This allows reuse of the T5 vocabulary and cross-entropy loss, making forecasting a language modeling problem over quantized tokens.
    \end{itemize}
    \end{exampleblock}
    }
    \hfill\quantlet{TSA\_ch11\_quiz2\_tokenization}{https://github.com/QuantLet/TSA/tree/main/TSA_ch11/TSA_ch11_quiz2_tokenization}
    \end{cminipage}
\end{frame}

\begin{frame}{Question 3}
    \begin{cminipage}{0.95\textwidth}
    \begin{alertblock}{Question}
        \begin{itemize}\setlength{\itemsep}{0pt}
            \item When should you prefer zero-shot inference over fine-tuning a foundation model?
        \end{itemize}
    \end{alertblock}

    \vspace{0.3cm}

    \begin{block}{Answer Choices}

        \textcolor{MainBlue}{\textbf{(A)}} When you have millions of labeled domain-specific observations\\[3pt]

        \textcolor{MainBlue}{\textbf{(B)}} When you need the highest possible accuracy on a single dataset\\[3pt]

        \textcolor{MainBlue}{\textbf{(C)}} When you have limited labeled data, need rapid deployment, or face cold-start problems\\[3pt]

        \textcolor{MainBlue}{\textbf{(D)}} When computational resources are not a concern

    \end{block}
    \end{cminipage}
\end{frame}

\begin{frame}{Question 3: Answer}
    \begin{cminipage}{0.95\textwidth}
    \vspace{-0.2cm}
    \begin{center}
        \includegraphics[width=0.98\textwidth, height=0.58\textheight, keepaspectratio]{ch11_quiz3_zeroshot_decision.pdf}
    \end{center}
    \vspace{-3mm}
    {\small
    \begin{exampleblock}{Answer: (C)}
    \begin{itemize}\setlength{\itemsep}{0pt}
        \item Zero-shot excels with limited data, cold-start, and rapid prototyping. The model's pre-trained knowledge transfers temporal patterns from billions of diverse series. Fine-tuning is preferred when sufficient domain data is available and marginal accuracy gains justify the cost.
    \end{itemize}
    \end{exampleblock}
    }
    \hfill\quantlet{TSA\_ch11\_quiz3\_zeroshot\_decision}{https://github.com/QuantLet/TSA/tree/main/TSA_ch11/TSA_ch11_quiz3_zeroshot_decision}
    \end{cminipage}
\end{frame}

\begin{frame}{Question 4}
    \begin{cminipage}{0.95\textwidth}
    \begin{alertblock}{Question}
        \begin{itemize}\setlength{\itemsep}{0pt}
            \item What is ``patching'' in the context of Transformer-based time series models, and why does it help?
        \end{itemize}
    \end{alertblock}

    \vspace{0.3cm}

    \begin{block}{Answer Choices}

        \textcolor{MainBlue}{\textbf{(A)}} Patching means filling missing values in the time series\\[3pt]

        \textcolor{MainBlue}{\textbf{(B)}} Patching groups $P$ consecutive points into one token, reducing attention complexity from $O(T^2)$ to $O((T/P)^2)$ while capturing local patterns\\[3pt]

        \textcolor{MainBlue}{\textbf{(C)}} Patching is a data augmentation technique that creates new training series\\[3pt]

        \textcolor{MainBlue}{\textbf{(D)}} Patching refers to updating model weights after deployment

    \end{block}
    \end{cminipage}
\end{frame}

\begin{frame}{Question 4: Answer}
    \begin{cminipage}{0.95\textwidth}
    \vspace{-0.2cm}
    \begin{center}
        \includegraphics[width=0.98\textwidth, height=0.58\textheight, keepaspectratio]{ch11_quiz4_patching.pdf}
    \end{center}
    \vspace{-3mm}
    {\small
    \begin{exampleblock}{Answer: (B)}
    \begin{itemize}\setlength{\itemsep}{0pt}
        \item Patching (PatchTST, 2023) groups $P$ consecutive points into a single token. This reduces the effective sequence length for attention, captures local temporal patterns within each patch, and enables the model to handle much longer context windows.
    \end{itemize}
    \end{exampleblock}
    }
    \hfill\quantlet{TSA\_ch11\_quiz4\_patching}{https://github.com/QuantLet/TSA/tree/main/TSA_ch11/TSA_ch11_quiz4_patching}
    \end{cminipage}
\end{frame}

\begin{frame}{Question 5}
    \begin{cminipage}{0.95\textwidth}
    \begin{alertblock}{Question}
        \begin{itemize}\setlength{\itemsep}{0pt}
            \item You need to forecast a multivariate energy dataset with 50 correlated variables at hourly frequency. Which foundation model is the best fit?
        \end{itemize}
    \end{alertblock}

    \vspace{0.3cm}

    \begin{block}{Answer Choices}

        \textcolor{MainBlue}{\textbf{(A)}} Chronos --- because it has the most model sizes to choose from\\[3pt]

        \textcolor{MainBlue}{\textbf{(B)}} TimeGPT --- because it is the easiest to use via API\\[3pt]

        \textcolor{MainBlue}{\textbf{(C)}} Moirai --- because it is the only model designed to handle any number of variates and any frequency natively\\[3pt]

        \textcolor{MainBlue}{\textbf{(D)}} Lag-Llama --- because it uses lag features that capture temporal dependencies

    \end{block}
    \end{cminipage}
\end{frame}

\begin{frame}{Question 5: Answer}
    \begin{cminipage}{0.95\textwidth}
    \vspace{-0.2cm}
    \begin{center}
        \includegraphics[width=0.98\textwidth, height=0.58\textheight, keepaspectratio]{ch11_quiz5_model_selection.pdf}
    \end{center}
    \vspace{-3mm}
    {\small
    \begin{exampleblock}{Answer: (C)}
    \begin{itemize}\setlength{\itemsep}{0pt}
        \item Moirai is the only foundation model designed for \textbf{any-variate} forecasting. It uses variate IDs and flexible patching to handle 50 correlated variables simultaneously. Chronos and Lag-Llama are univariate; they would process each variable independently, ignoring cross-variable correlations.
    \end{itemize}
    \end{exampleblock}
    }
    \hfill\quantlet{TSA\_ch11\_quiz5\_model\_selection}{https://github.com/QuantLet/TSA/tree/main/TSA_ch11/TSA_ch11_quiz5_model_selection}
    \end{cminipage}
\end{frame}

%=============================================================================
% BIBLIOGRAPHY
%=============================================================================
\begin{frame}{Bibliography I}
    \begin{block}{Foundational Papers}
        {\footnotesize
        \begin{itemize}\setlength{\itemsep}{0pt}
            \item Vaswani, A., Shazeer, N., Parmar, N., et al.\ (2017). Attention Is All You Need, \textit{Advances in Neural Information Processing Systems (NeurIPS)}, 30, 5998--6008.
            \item Ansari, A.F., Stella, L., Turkmen, C., et al.\ (2024). Chronos: Learning the Language of Time Series, \textit{arXiv preprint arXiv:2403.07815}.
            \item Das, A., Kong, W., Leber, A., et al.\ (2024). A Decoder-Only Foundation Model for Time-Series Forecasting, \textit{Proceedings of ICML 2024}.
        \end{itemize}
        }
    \end{block}

    \begin{exampleblock}{Foundation Models}
        {\footnotesize
        \begin{itemize}\setlength{\itemsep}{0pt}
            \item Garza, A., \& Mergenthaler-Canseco, M.\ (2024). TimeGPT-1, \textit{International Journal of Forecasting}, forthcoming.
            \item Rasul, K., Ashok, A., Williams, A.R., et al.\ (2024). Lag-Llama: Towards Foundation Models for Probabilistic Time Series Forecasting, \textit{Proceedings of NeurIPS 2024}.
        \end{itemize}
        }
    \end{exampleblock}
\end{frame}

\begin{frame}{Bibliography II}
    \begin{block}{Architecture and Pre-Training}
        {\footnotesize
        \begin{itemize}\setlength{\itemsep}{0pt}
            \item Woo, G., Liu, C., Kumar, A., et al.\ (2024). Unified Training of Universal Time Series Forecasting Transformers, \textit{Proceedings of ICML 2024}.
            \item Nie, Y., Nguyen, N.H., Sinthong, P., \& Kalagnanam, J.\ (2023). A Time Series is Worth 64 Words: Long-term Forecasting with Transformers, \textit{Proceedings of ICLR 2023}.
            \item Bommasani, R., Hudson, D.A., Adeli, E., et al.\ (2021). On the Opportunities and Risks of Foundation Models, \textit{arXiv preprint arXiv:2108.07258}.
        \end{itemize}
        }
    \end{block}

    \begin{alertblock}{Benchmarks and Evaluation}
        {\footnotesize
        \begin{itemize}\setlength{\itemsep}{0pt}
            \item Godahewa, R., Bergmeir, C., Webb, G., et al.\ (2021). Monash Time Series Forecasting Archive, \textit{Proceedings of NeurIPS 2021 Datasets and Benchmarks Track}.
            \item Hu, E.J., Shen, Y., Wallis, P., et al.\ (2022). LoRA: Low-Rank Adaptation of Large Language Models, \textit{Proceedings of ICLR 2022}.
        \end{itemize}
        }
    \end{alertblock}
\end{frame}

\begin{frame}{Bibliography III}
    \begin{block}{Additional References}
        {\footnotesize
        \begin{itemize}\setlength{\itemsep}{0pt}
            \item Kaplan, J., McCandlish, S., Henighan, T., et al.\ (2020). Scaling Laws for Neural Language Models, \textit{arXiv preprint arXiv:2001.08361}.
            \item Makridakis, S., Spiliotis, E., \& Assimakopoulos, V.\ (2018). The M4 Competition: 100,000 Time Series and 61 Forecasting Methods, \textit{International Journal of Forecasting}, 36(1), 54--74.
            \item Petropoulos, F., et al.\ (2022). Forecasting: Theory and Practice, \textit{International Journal of Forecasting}, 38(3), 845--1054.
        \end{itemize}
        }
    \end{block}
    \vspace{-1mm}
    \begin{exampleblock}{Online Resources and Code}
        {\footnotesize
        \begin{itemize}\setlength{\itemsep}{0pt}
            \item \textbf{Quantlet}: \url{https://quantlet.com} -- Code platform for quantitative methods
            \item \textbf{Quantinar}: \url{https://quantinar.com} -- Learning platform for quantitative methods
            \item \textbf{GitHub TSA}: \url{https://github.com/QuantLet/TSA/tree/main/TSA_ch11} -- Python code for this chapter
        \end{itemize}
        }
    \end{exampleblock}
\end{frame}

\begin{frame}{}
    \centering
    \Huge\textcolor{MainBlue}{Thank You!}

    \vspace{1cm}

    \Large Questions?

    \vspace{0.8cm}

    \normalsize

    Course materials available at: \url{https://danpele.github.io/Time-Series-Analysis/}

    \vspace{0.2cm}

    \href{https://quantlet.com}{\raisebox{-0.15em}{
\includegraphics[height=0.8em]{ql_logo.png}} Quantlet} \hspace{0.5cm}
    \href{https://quantinar.com}{\raisebox{-0.15em}{
\includegraphics[height=0.8em]{qr_logo.png}} Quantinar}
\end{frame}

\end{document}
